\begin{abstract}
Repair guidance for a damaged garment requires a reliable observation and a procedure compatible with the fabric and available tools. We present UCRA, an offline Android workflow combining image-quality screening, selective hole localization, user-confirmed context, source-linked procedure retrieval, and a repair journal. We built an exploratory benchmark from the CISUTAC garment-hole export with 407 views of 406 garments, including 61 held-out garments. We compared a logistic baseline, a random forest, and whole-view and tiled center-heatmap detectors. At validation-selected operating points, the random forest localized 3 of 61 views at intersection over union (IoU) of 0.25 while showing 53 candidates. The tiled detector localized 8 of 61 views while showing 18 candidates, with 7 reaching IoU 0.50. Its score threshold increased selective precision from 19.7\% to 44.4\% while reducing correct proposals from 12 to 8. All evaluated views contained annotated holes, leaving image-level specificity unknown. The results support using the detector as an experimental aid that requires user confirmation. Reliable handheld detection and procedure compatibility require external negative controls and independent evaluation.
\end{abstract}
